\section{ПРЕДМЕТНАЯ ОБЛАСТЬ И ТЕХНИЧЕСКОЕ ЗАДАНИЕ}

\subsection{Точное животноводство и автоматизация управления молочными фермами}

Концепция точного животноводства (precision livestock farming, PLF)~\cite{PrecisionLivestock} предполагает непрерывный автоматизированный мониторинг состояния животных с использованием датчиков и информационных технологий для принятия управленческих решений в реальном времени. Данная концепция была предложена Д.~Беркмансом и К.~Беркманс в 2008~году и с тех пор получила широкое распространение в молочном животноводстве~\cite{DairyPrecision}. Согласно исследованию Европейской комиссии по сельскому хозяйству, внедрение PLF-технологий позволяет снизить использование антибиотиков на 20--30~процентов, повысить продуктивность стада на 10--15~процентов и сократить трудозатраты на мониторинг в 3--5~раз.

В молочном животноводстве PLF охватывает несколько ключевых направлений:
\begin{enumerate}
    \item мониторинг надоев --- объём молока по четвертям вымени, скорость доения (кг/мин), электропроводность (мСм/см) как индикатор мастита, количество соматических клеток (\acrshort{scc}, тыс./мл). Типичная корова при доильном роботе Lely Astronaut генерирует 2--3~визита в сутки, каждый из которых содержит 15--20 параметров;
    \item мониторинг здоровья --- активность (шаги за час), жвачка (минуты жевания за период), температура тела и индекс потребления корма. Датчики активности (Lely Qwes, SCR HR) фиксируют данные с дискретностью 1--2~часа, что позволяет обнаруживать отклонения в течение суток;
    \item управление воспроизводством --- обнаружение охоты по активности (повышение в 3--5~раз относительно базовой линии), контроль стельности и планирование отёлов и запусков. Интервал между отёлами --- ключевой KPI: целевое значение 365--400~дней, каждый лишний день обходится в 3--5~долларов потерь от непродуктивного содержания;
    \item управление кормлением --- контроль рационов (сухое вещество, netto-energy-lactation), потребление кормов (кг/день) и остатки; эффективность корма (л молока / кг сухого вещества) является одним из главных экономических показателей;
    \item управление оборудованием --- отслеживание состояния доильных роботов, аномалий работы и прогнозирование технического обслуживания.
\end{enumerate}

Роботизированные доильные установки, такие как Lely Astronaut, автоматически собирают данные о каждом доении. Датчики активности (Lely Qwes) фиксируют шаги, жвачку и время приёма корма. Эти данные передаются в облачный сервис Lely Horizon через REST API. Совокупный объём данных для стада из 500~коров составляет до 250\,000 записей в сутки, что делает невозможной их ручную обработку.

Раннее выявление таких заболеваний, как мастит и кетоз, критически важно для экономической эффективности. По данным международного исследования~\cite{MastitisDetection}, средняя стоимость случая клинического мастита составляет от 100 до 300~долларов США с учётом снижения надоев (70~процентов потерь), затрат на лечение (15~процентов) и преждевременной выбраковки (15~процентов). Кетоз поражает до 40~процентов коров в ранний послеродовой период и при невыявлении приводит к потере до 500~кг молока за лактацию. Модели машинного обучения способны анализировать паттерны в многомерных временных рядах и обнаруживать отклонения на ранних стадиях, когда клинические признаки ещё не очевидны~\cite{MastitisDetection}.

\subsection{Обзор существующих решений}

На рынке представлены несколько информационных систем управления молочными фермами. Сравнительный анализ приведён в таблице~\ref{solcompare}.

\begin{table}
  \caption{Сравнительный анализ существующих решений}
  \small
  \begin{tabularx}{\textwidth}{|p{3.3cm}|X|X|X|X|}
    \hline
    \textbf{Критерий} & \textbf{Lely Horizon} & \textbf{DeLaval DelPro} & \textbf{SCR Heatime} & \textbf{1С: Управление фермой} \\
    \hline
    Учёт животных & Да & Да & Нет & Да \\
    \hline
    Учёт надоев & Да & Да & Нет & Да \\
    \hline
    Воспроизводство & Да & Да & Да & Да \\
    \hline
    Кормление & Да & Да & Нет & Ограничено \\
    \hline
    Интеграция с роботами & Lely & DeLaval & Нет & Нет \\
    \hline
    Модели ML & Нет & Нет & Да (охота) & Нет \\
    \hline
    Открытый API & Ограничен & Нет & Нет & Нет \\
    \hline
    Кастомизация & Низкая & Низкая & Низкая & Средняя \\
    \hline
    Стоимость & Высокая & Высокая & Средняя & Низкая \\
    \hline
  \end{tabularx}
  \label{solcompare}
\end{table}

\textbf{Lely Horizon} --- облачная платформа от производителя доильных роботов Lely. Обеспечивает мониторинг стада и оборудования в реальном времени: отображение текущего надоя, активности, жвачки и состояния здоровья по каждой корове. Интерфейс поддерживает настройку пороговых оповещений (например, падение надоя > 20~процентов за сутки). Однако система ограничена экосистемой Lely: не поддерживает оборудование других производителей, не предоставляет API для интеграции сторонних моделей аналитики и не позволяет добавлять собственные правила оповещений. Для ферм, использующих смешанное оборудование, это является критическим ограничением.

\textbf{DeLaval DelPro} --- система управления фермой от DeLaval. Предлагает учёт животных и надоев, управление доильным залом и мониторинг оборудования. Является проприетарной и привязана к оборудованию DeLaval: данные доильных роботов других производителей недоступны. Не поддерживает интеграцию с оборудованием сторонних производителей, что затрудняет внедрение на фермах с гетерогенным оборудованием. Предиктивная аналитика ограничена базовыми пороговыми оповещениями.

\textbf{SCR Heatime} --- система мониторинга активности и обнаружения охоты от SCR Engineers (ныне Allflex). Система специализируется на воспроизводстве: использует данные о активности (шаги) и жвачке для обнаружения охоты с помощью запатентованного алгоритма. Встроенная модель машинного обучения обеспечивает чувствительность обнаружения охоты до 95~процентов при специфичности 90~процентов~\cite{PrecisionLivestock}. Однако SCR Heatime не охватывает полный цикл управления фермой: нет учёта кормления, надоев, запасов и оборудования. Система является закрытой и не предоставляет API.

\textbf{1С: Управление фермой} --- отечественное решение на платформе 1С:Предприятие. Обеспечивает базовый учёт животных, надоев, воспроизводства и кормления. Преимуществом является локализация, соответствие российскому законодательству и низкая стоимость. Однако система не поддерживает интеграцию с роботизированным оборудованием, не имеет модулей машинного обучения, ограничена в кастомизации отчётов и использует монолитную архитектуру, затрудняющую масштабирование.

Анализ показывает, что ни одно из существующих решений не обеспечивает одновременно:
\begin{enumerate}
    \item полный цикл управления стадом от учёта до аналитики;
    \item интеграцию с оборудованием сторонних производителей через открытое API;
    \item встроенные модели машинного обучения для предиктивной аналитики;
    \item открытую расширяемую архитектуру для добавления собственных моделей.
\end{enumerate}

\subsection{Техническое задание}

На основании проведённого анализа сформулируем техническое задание на разработку информационной системы.

Учёт животных должен обеспечивать ведение реестра с полной биографической информацией (инвентарный номер, кличка, пол, дата рождения, порода, происхождение), историю переводов между группами и местоположениями, деактивацию без удаления записей для сохранения истории, поиск по кличке и инвентарному номеру с поддержкой нечёткого поиска (триграммное сходство pg\_trgm), импорт из CSV и экспорт карточки в PDF, а также пакетные операции для массовой деактивации и обновления до 500~животных.

Учёт надоев молока должен включать суточные надои (объём в кг, жир, белок, лактозу, SCC, FPR), визиты на доение (время начала и окончания, объём по четвертям вымени, скорость доения, данные робота), результаты лабораторных анализов качества молока и мониторинг резервуара-охладителя.

Управление воспроизводством должно поддерживать регистрацию отёлов с записью данных о телёнке, осеменений с указанием быка-производителя, контроль стельности, регистрацию охот на основе данных активности и управление запусками (сухостойный период).

Управление кормлением должно обеспечивать работу с типами кормов и группами кормов, дневными объёмами (плановое и фактическое потребление) и визитами на кормление (время, объём, остатки).

Фитнес-мониторинг и пастьба должны охватывать активность (шаги за период), жвачку (минуты жевания за период) и данные о пастьбе.

Отчёты и аналитика должны включать генерацию отчётов более 17 типов с экспортом в CSV и PDF, ключевые показатели эффективности (KPI), тренды надоев с прогнозом и прогноз стельных коров.

Интеграция с роботизированным оборудованием должна обеспечивать автоматическую синхронизацию с облачным API Lely Horizon с настраиваемым интервалом для 23 типов данных и шифрованием учётных данных (AES-256-GCM).

Предиктивная аналитика должна включать:
\begin{enumerate}
    \item прогноз надоев на корову (с доверительными интервалами q10/q90) и на стадо;
    \item обнаружение мастита и кетоза;
    \item выявление охоты;
    \item оценку риска выбраковки;
    \item кластеризацию стада (KMeans);
    \item обнаружение аномалий оборудования (Isolation Forest);
    \item рекомендации по кормлению;
    \item индекс здоровья на основе Z-оценок;
    \item сравнение моделей прогнозирования (8 моделей с автоматическим выбором).
\end{enumerate}

Нефункциональные требования включают:
\begin{enumerate}
    \item аутентификацию и авторизацию с разграничением ролей (admin/user);
    \item адаптивный веб-интерфейс для настольных и мобильных устройств;
    \item контейнеризацию всех компонентов (Docker);
    \item мониторинг работоспособности (Prometheus, Grafana, Loki, Tempo);
    \item автоматизированное развёртывание через CI/CD (GitHub Actions).
\end{enumerate}

\subsection{Стек технологий}

Выбор технологий обоснован требованиями к производительности, безопасности и удобству разработки.

Для серверной части выбран язык программирования Rust~\cite{RustLang}, обеспечивающий безопасность работы с памятью без сборщика мусора за счёт системы владения (ownership) и заимствования (borrowing), что исключает целый класс уязвимостей: использование после освобождения (use-after-free), двойное освобождение (double-free), выход за границы массива. По результатам бенчмарков TechEmpower Framework Benchmarks, веб-фреймворк Actix-web (аналогичного уровня с Axum) демонстрирует пропускную способность более 300\,000 запросов/сек на одном ядре, что в 3--5 раз выше Node.js/Express и сопоставимо с C++. Веб-фреймворк Axum~\cite{Axum} построен на асинхронном рантайме Tokio~\cite{Tokio} и библиотеке Tower~\cite{Tower}, что обеспечивает высокую пропускную способность при обработке HTTP-запросов за счёт неблокирующего ввода-вывода. Библиотека SQLx~\cite{SQLx} выполняет проверку SQL-запросов при компиляции (compile-time query checking), исключая ошибки в SQL-коде на этапе сборки.

В качестве основной СУБД выбран PostgreSQL~\cite{PostgreSQL} --- зрелая (более 25 лет разработки), надёжная (ACID-транзакции) реляционная база данных с расширенной поддержкой типов данных, индексов и хранимых процедур. Расширение TimescaleDB~\cite{TimescaleDB} обеспечивает эффективное хранение и агрегацию временных рядов с использованием гипертаблиц~\cite{TimescaleHypertable} и непрерывных агрегатов, что критично для таблиц с суточными надоями, активностью и жвачкой (миллионы записей за год).

Для клиентского приложения выбран фреймворк SvelteKit~\cite{SvelteKit}. В отличие от React и Angular, Svelte~5~\cite{Svelte5} не использует виртуальный \acrshort{dom}, а компилирует компоненты в оптимизированный императивный JavaScript-код: размер бандла в 2--3 раза меньше React-аналога, а время рендеринга --- в 1,5--2 раза быстрее. SvelteKit поддерживает серверный рендеринг (\acrshort{ssr}), улучшающий скорость начальной загрузки. Tailwind CSS~\cite{TailwindCSS} используется для стилизации (utility-first подход с автоматическим удалением неиспользуемых классов), Chart.js~\cite{ChartJS} --- для визуализации данных.

Для модуля машинного обучения выбран Python как основной язык в области data science с наибольшей экосистемой библиотек. FastAPI~\cite{FastAPI} обеспечивает высокопроизводительный API с автоматической генерацией документации (OpenAPI/Swagger) и валидацией типов через Pydantic. XGBoost~\cite{XGBoost} используется для задач классификации и регрессии благодаря высокой производительности (в 10--20 раз быстрее scikit-learn на тех же данных) и поддержке квантильной регрессии. Prophet~\cite{Prophet} --- для прогнозирования на уровне стада с автоматической декомпозицией тренда и сезонности. Optuna~\cite{Optuna} --- для байесовской оптимизации гиперпараметров (TPE-сэмплер).

Инфраструктура базируется на Docker~\cite{Docker} и Docker Compose~\cite{DockerCompose}, обеспечивающих контейнеризацию и оркестрацию с изолированными средами для каждого компонента. Nginx~\cite{Nginx} выступает в роли обратного прокси с терминацией SSL, rate limiting и кэшированием статики. Стек мониторинга Grafana LGTM (Prometheus~\cite{Prometheus}, Grafana~\cite{Grafana}, Loki~\cite{Loki}, Tempo~\cite{Tempo}) обеспечивает комплексное наблюдение: метрики, логи и распределённые трассировки. GitHub Actions~\cite{GHActions} автоматизирует CI/CD: параллельное тестирование трёх компонентов, сборку Docker-образов и деплой.
